\documentclass[tikz, border=6pt]{standalone}

\usepackage{xeCJK}
\setCJKmainfont{Microsoft JhengHei}
\setCJKsansfont{Microsoft JhengHei}
\setmainfont{Arial}
\setsansfont{Arial}
\usepackage{amsmath, amssymb}

\usetikzlibrary{positioning, arrows.meta, shapes.geometric, fit, backgrounds, calc, shadows}

\begin{document}

\begin{tikzpicture}[
    node distance=1.1cm and 1.5cm,
    % 基礎模組方塊
    module/.style={rectangle, rounded corners=8pt, draw=gray!60, thick, fill=white, minimum width=6.2cm, minimum height=1.7cm, align=center, font=\sffamily\large, drop shadow={opacity=0.08, shadow xshift=2pt, shadow yshift=-2pt}},
    highlight/.style={module, draw=orange!60, fill=orange!5},
    % 節點樣式
    user_node/.style={circle, draw=blue!60, fill=blue!10, inner sep=5pt, font=\sffamily\large\bfseries},
    item_node/.style={circle, draw=orange!60, fill=orange!10, inner sep=5pt, font=\sffamily\large\bfseries},
    entity_node/.style={diamond, draw=green!60, fill=green!10, inner sep=4pt, font=\sffamily\large\bfseries},
    % 箭頭樣式
    arrow/.style={-{Stealth[scale=1.4]}, thick, draw=gray!70},
    dashedarrow/.style={-{Stealth[scale=1.4]}, thick, dashed, draw=gray!70},
    % 外框背景
    framebox/.style={rectangle, rounded corners=12pt, draw=gray!40, dashed, thick, inner sep=18pt}
]

% 色彩定義
\definecolor{slate}{RGB}{112, 128, 144}
\definecolor{bg_input}{RGB}{245, 247, 250}
\definecolor{bg_prop}{RGB}{253, 248, 243}
\definecolor{bg_pred}{RGB}{243, 250, 246}

% ==========================================
% 1. 輸入層 (Col 1: Center x = 0.0)
% ==========================================
\node[user_node] (u1) at (-2.8, 2.8) {用};
\node[item_node] (i1) at (0.0, 2.8) {食};
\node[entity_node] (e1) at (2.8, 2.8) {實};

% interact 顯著上移 (yshift=10pt)，has_ingredient 顯著下移 (yshift=-12pt)，確保 100% 無接觸
\draw[arrow] (u1) -- node[above, font=\sffamily\small\bfseries, yshift=10pt] {interact} (i1);
\draw[arrow] (i1) -- node[below, font=\sffamily\small\bfseries, yshift=-12pt] {has\_ingredient} (e1);

\node[module, minimum width=6.2cm] (init) at (0.0, 0.0) {\Large\bfseries 嵌入向量初始化\\[0.1cm]\large (使用者、食譜、實體 $\rightarrow$ 向量 $\mathbf{e}^{(0)}$)};

% ==========================================
% 2. 傳播層 (Col 2: Center x = 9.0，拉開大間距)
% ==========================================
\node[highlight, minimum width=6.2cm] (attn) at (9.0, 2.8) {\Large\bfseries 關係感知注意力機制\\[0.1cm]\large (計算動態權重 $\pi(h,r,t)$)};
\node[module, minimum width=6.2cm] (agg) at (9.0, 0.0) {\Large\bfseries 雙向互動聚合器 (Bi-Interaction)\\[0.1cm]\large (特徵組合 $\mathbf{e}_h^{(l)}$)};

\draw[arrow] (attn) -- (agg);
\draw[arrow, bend right=55] (agg.east) to node[right=5pt, align=left, font=\sffamily\small\bfseries] {重複傳播\\$L$ 層 ($L=3$)} (attn.east);

% ==========================================
% 3. 預測層 (Col 3: Center x = 18.0，拉開大間距)
% ==========================================
\node[module, minimum width=6.2cm] (concat) at (18.0, 2.8) {\Large\bfseries 跨層特徵拼接 (JK Net)\\[0.1cm]\large (拼接特徵 $\mathbf{e}^* = \mathbf{e}^{(0)} \| \dots \| \mathbf{e}^{(L)}$)};
\node[module, minimum width=6.2cm] (pred) at (18.0, 0.0) {\Large\bfseries 內積運算與偏好預測\\[0.1cm]\large (預測分數 $\hat{y}(u,i) = {\mathbf{e}_u^*}^T \mathbf{e}_i^*$)};

\draw[arrow] (concat) -- (pred);

% ==========================================
% 外框背景 ( 用 Dummy 點強制 3 個外框 100% 相同頂底高度 )
% ==========================================
\begin{scope}[on background layer]
    \coordinate (T1) at (0.0, 4.0);
    \coordinate (B1) at (0.0, -0.9);
    \node[framebox, fill=bg_input, fit=(u1)(e1)(init)(T1)(B1)] (input_frame) {};
    \node[anchor=north west, font=\sffamily\bfseries\Large, text=slate, xshift=8pt, yshift=-6pt] at (input_frame.north west) {1. 輸入層：協同知識圖 (CKG)};

    \coordinate (T2) at (9.0, 4.0);
    \coordinate (B2) at (9.0, -0.9);
    \node[framebox, fill=bg_prop, fit=(attn)(agg)(T2)(B2)] (prop_frame) {};
    \node[anchor=north west, font=\sffamily\bfseries\Large, text=slate, xshift=8pt, yshift=-6pt] at (prop_frame.north west) {2. 注意力訊息傳播層};

    \coordinate (T3) at (18.0, 4.0);
    \coordinate (B3) at (18.0, -0.9);
    \node[framebox, fill=bg_pred, fit=(concat)(pred)(T3)(B3)] (pred_frame) {};
    \node[anchor=north west, font=\sffamily\bfseries\Large, text=slate, xshift=8pt, yshift=-6pt] at (pred_frame.north west) {3. 預測層};
\end{scope}

% ==========================================
% 跨欄連接箭頭
% ==========================================
\draw[arrow, thick] (init.east) -- node[above, font=\sffamily\normalsize\bfseries] {初始特徵} (agg.west |- init.east);
\draw[arrow, thick] (agg.east) -- node[above, font=\sffamily\normalsize\bfseries] {高階特徵} (pred.west |- agg.east);

% Hop 0 初始嵌入長虛線：從 Column 1 框底往下繞過 3 個外框下方，向右上指入 pred_frame 底部
\draw[dashedarrow] (input_frame.south) -- ++(0, -0.6) -| node[pos=0.25, below, font=\sffamily\Large\bfseries] {初始嵌入 (Hop 0)} (pred_frame.south);

\end{tikzpicture}

\end{document}